% !PRINTED 289
\weberpage{289}
%<<CONT
\[
\begin{array}{ll}
1. & \tfrac{1}{2}\big[E(f; \varphi, \psi) - A(f; \varphi, \psi)\big], \\
2. & \tfrac{1}{2}\big[E(\varphi; \psi, f) - A(\varphi; \psi, f)\big], \\
3. & \tfrac{1}{2}\big[E(\psi; f, \varphi) - A(\psi; f, \varphi)\big], \\
4. & \tfrac{1}{2}\big[E(f; \psi, \varphi) - A(f; \psi, \varphi)\big], \\
5. & \tfrac{1}{2}\big[E(\psi; \varphi, f) - A(\psi; \varphi, f)\big], \\
6. & \tfrac{1}{2}\big[E(\varphi; f, \psi) - A(\varphi; f, \psi)\big],
\end{array}
\]
\gesperrt{and there now holds the fundamental theorem that these six determinations
yield the same number.}

To prove it, we traverse the curve $f$ in the positive
sense, and attend to all the intersection points with $\varphi$
and $\psi$. The path along the curve $f$ must in doing so enter
the inner region $(\varphi, \psi)$ just as often as it steps out of it,
because it must return again to its starting point.
But the curve $f$ enters at the points $E(f; \varphi, \psi)$
and $A(f; \psi, \varphi)$ and steps out at the points $A(f; \varphi, \psi)$ and
$E(f; \psi, \varphi)$. Hence
\[
E(f; \varphi, \psi) + A(f; \psi, \varphi) = A(f; \varphi, \psi) + E(f; \psi, \varphi),
\]
by which the agreement of 1. and 4. is established.

\weberfigureimage{weber-b1-fig09.png}{0.42}

If now, secondly, we traverse the curve $\varphi$ in the negative sense,
and attend to its intersection points with $f$, it turns out
that every point $E(f; \varphi, \psi)$ is at the same time a point $E(\varphi; f, \psi)$,
and every point $A(f; \varphi, \psi)$ a point $A(\varphi; f, \psi)$. For in
Fig.~9, in which the point $\alpha$ represents any one of the intersection points
of $f$ and $\varphi$, this point is an $E(f; \varphi, \psi)$ and
an $E(\varphi; f, \psi)$ when $\psi$ is positive at $\alpha$, and an $A(f; \varphi, \psi)$ and
an $A(\varphi; f, \psi)$ when $\psi$ is negative at $\alpha$.

From this follows the agreement of 1. with 6. There follows
further, by a renewed application of the first argument, the
agreement of 6. with 2., and then by application of the
second argument that of 2. with 5., and finally by the first that
of 5. with 3., \gesperrt{hence the agreement of all six expressions.}

One is therefore justified in designating the number so determined
simply as the \gesperrt{characteristic of the system of functions} $(f, \varphi, \psi)$.
%CONT>>
